\section{Stable certificates for exact resource profiles}
\label{sec:v12-certificates}

The task-cover theorem and the analytic comparison theorem concern the
same invariant: the least common-encoder excess at a prescribed width
profile. The following consequence states precisely when an approximate
experiment certifies an exact resource answer. This is also the point at
which the microscopic decision limit can be used without exchanging a
model limit with optimization over encoders.

For the finite restart system of Section~\ref{sec:v12-support}, write
\[
 V_E(S,q)=\inf_\pi\max_{t,d}
  \left[r^E_{td}(\pi)-\sum_h\eta_t(h)m^E_{td}(h)\right].
\]
If a positive entry regret exists, define $g_E=\alpha\gamma$ as in
Corollary~\ref{cor:v12-gap}; for a Brier task family use
$g_E=\alpha\delta^2/2$. This gap is defined by the target task table, not
estimated from a lossy experiment. In the trivial case of no positive
regret every resource profile is feasible and no gap test is necessary.

\begin{theorem}[Comparison, completion, and certification]
\label{thm:v12-certificate}
Let $E$ be a finite restart experiment with $g_E>0$, and let $F$ be an
experiment with the same decision rows and loss scale. Suppose two-sided
cost-one comparison at the declared seed signature gives risk-body
Hausdorff error at most $\varepsilon$, and the corresponding full-history
benchmarks differ by at most $\varepsilon$. Then
\begin{equation}
 |V_F(S,q)-V_E(S,q)|\le2\varepsilon.
 \label{eq:v12-certificate-error}
\end{equation}
If $\varepsilon<g_E/4$, the single test
\begin{equation}
 V_F(S,q)<g_E/2
 \label{eq:v12-test}
\end{equation}
is equivalent to existence of a causally closed compatible cover for $E$
with width profile $S$. Equality of the exact feasible profiles holds for
all private and independent public seed signatures, even though the
positive-excess values for those signatures may differ.

More generally, suppose only a simulator from $E$ to $F$ is given, with
cost $(K,r)$, and suppose its risk and benchmark errors are at most
$\varepsilon$. If
\[
 V_F(S,q)+2\varepsilon<g_E,
\]
then $E$ has a causally closed compatible cover of width at most $KS$.
The public seed bound of the implemented design is $rq$, but it is
unnecessary for exactness of that cover.
\end{theorem}

\begin{proof}
The two risk-body inclusions bound the infimum of the maximum centered
row by the opposite value plus the risk error and the benchmark error.
This proves \eqref{eq:v12-certificate-error}. By
Corollary~\ref{cor:v12-gap}, $V_E(S,q)$ is either zero or at least $g_E$,
and it is zero precisely for a causally closed compatible cover. In the
first case $V_F(S,q)\le2\varepsilon<g_E/2$; in the second case
$V_F(S,q)\ge g_E-2\varepsilon>g_E/2$. This proves the test without
requiring an optimizer for $F$. The one-sided implementation gives
$V_E(KS,rq)\le V_F(S,q)+2\varepsilon<g_E$, forcing its value to be zero.
Theorem~\ref{thm:v12-support} then supplies a deterministic cover machine
at width $KS$.
\end{proof}

\begin{corollary}[Certified approximation sequence]
\label{cor:v12-cert-sequence}
For any family $F_n$ satisfying the hypotheses of the preceding theorem
with a known error bound $\varepsilon_n\to0$, every fixed resource
profile of $E$ is classified correctly by \eqref{eq:v12-test} once
$\varepsilon_n<g_E/4$. If the computed value has an additional certified
numerical error $\kappa_n$, a sufficient condition is
$2\varepsilon_n+\kappa_n<g_E/2$, with the corresponding interval test.
In applications of Theorem~\ref{thm:v12-physical}, the error bound is
$E_n$ whenever the entire declared restart family, including its
preparations, is among the compared physical experiments.
\end{corollary}

\begin{proof}
Apply \eqref{eq:v12-certificate-error} and add the numerical error.
The last qualification ensures that the physical comparison controls
these restart rows, rather than a different initial preparation or only
a single uncontrolled experiment.
\end{proof}

The gap may be small, and computing the finite common-encoder optimum
can be difficult. The statement is a certification theorem, not a
polynomial-time algorithm. In particular it does not make a constrained
inverse randomization assertion: a risk witness need not manufacture a
simulator of a prescribed cost. Its content is the stability of exact
causal task completion under an independently verified, resource-priced
comparison.
